\par\medskip\noindent\textbf{Cómo leer la salida.}
{\small\setlength{\tabcolsep}{4pt}\renewcommand{\arraystretch}{1.15}
\begin{xltabular}{\linewidth}{@{}>{\raggedright\arraybackslash}p{0.24\linewidth}>{\raggedleft\arraybackslash}p{0.15\linewidth}>{\raggedright\arraybackslash}X@{}}
\toprule Indicador & Valor & Interpretación \\ \midrule \endfirsthead
\toprule Indicador & Valor & Interpretación \\ \midrule \endhead
\bottomrule \endlastfoot
Días con rendimiento cero & 171 & Se excluyen porque no tienen signo. \\
\texttt{n1(+)}, \texttt{n2(-)} & p. ej. 4617 / 4333 & Número de rendimientos positivos y negativos. \\
\texttt{rachas} & 4395 / 343 / 383 & Secuencias consecutivas de igual signo observadas. \\
\texttt{esperadas} & 4471,5 / 393,1 / 385,3 & Rachas esperadas si los signos fueran independientes, $1+2n_1n_2/n$. \\
\texttt{z} Brent & $-1{,}62$ & Menos rachas de las esperadas, pero no significativo: no hay dependencia en los signos. \\
\texttt{z} S1V3 & $-3{,}66$ & Muchas menos rachas: los signos de la cartera pequeña se repiten (persistencia), coherente con su $\rho_1$ y $VR>1$. \\
\texttt{z} S5V3 & $-0{,}17$ & La cartera grande no muestra dependencia, igual que en la razón de varianzas. \\
\end{xltabular}}
